% Problem record op_c1499a0073ee800c. This TeX file is derived from
% database/problems_json/op_c1499a0073ee800c.json by scripts/sync-tex.mjs; edit the JSON record
% and rerun the script rather than editing this file.
\section{Additivity of entanglement of formation for two-mode Gaussian states}
\label{sec:op_c1499a0073ee800c}

\paragraph{Problem.}
Is unrestricted entanglement of formation additive on tensor powers of every finite-energy two-mode Gaussian state?

Let $\rho_{AB}$ be any finite-energy two-mode Gaussian density operator, shared between one mode at each party. Finite energy means finite total mean photon number. Define the unrestricted pure-state convex roof

\begin{equation}
\begin{gathered}
E_F(\omega):=\inf_\mu\int S(\operatorname{Tr}_B|\psi\rangle\langle\psi|)\,\mu(d\psi),
\\
S(\tau):=-\operatorname{Tr}(\tau\log_2\tau),
\end{gathered}
\label{eq:ser12-statement-1}
\end{equation}

In Eq.~\eqref{eq:ser12-statement-1}, $\mu$ ranges over pure-state probability measures with barycentre $\omega$.

The target is $E_F(\rho_{AB}^{\otimes k})=kE_F(\rho_{AB})$ for every integer $k\geq1$. The bipartition for the tensor power is $A^k:B^k$. Non-Gaussian vectors are allowed in all decompositions.

\subsection*{Status}

Unsolved

\subsection*{Source}

Adesso explicitly leaves equality of entanglement cost and single-copy entanglement of formation open in the Conclusion and Outlook \sourcecite{ref:ser12-2}{Ade26}. Under the finite-entropy cost theorem this is equivalent to tensor-power additivity \sourcecite{ref:ser12-1}{YKHL25}.

\subsection*{Progress}

\begin{itemize}
  \item Theorem 7 of Yamasaki et al. identifies entanglement cost with the regularized convex roof:

\begin{equation}
E_C(\rho_{AB})=\lim_{k\to\infty}\frac1kE_F(\rho_{AB}^{\otimes k}),
\label{eq:ser12-progress-1-1}
\end{equation}

Equation~\eqref{eq:ser12-progress-1-1} includes the continuous pure-state roof. Finite local entropy suffices; finite-energy states on finitely many modes satisfy this hypothesis. Here $E_C$ is the asymptotic Bell-pair consumption per target copy under LOCC with vanishing trace-norm error. \sourcecite{ref:ser12-1}{YKHL25}

  \item Tensor products of single-copy ensembles give only

\begin{equation}
E_C(\rho_{AB})\leq\frac1kE_F(\rho_{AB}^{\otimes k})\leq E_F(\rho_{AB}),
\label{eq:ser12-progress-2-1}
\end{equation}

The reverse inequality needed to saturate Eq.~\eqref{eq:ser12-progress-2-1} remains open, as stated in Adesso's Conclusion and Outlook. \sourcecite{ref:ser12-2}{Ade26}

  \item Wilde, Proposition 6 and Eqs. (141)--(143), proves Eq.~\eqref{eq:ser12-progress-3-1} for $\rho_{AB}=(\operatorname{id}\otimes\mathcal N)(|\Phi_{N_s}\rangle\langle\Phi_{N_s}|)$. Here $|\Phi_{N_s}\rangle:=\sum_{j\geq0}\sqrt{N_s^j/(N_s+1)^{j+1}}|j,j\rangle$ is a two-mode squeezed vacuum with $N_s\geq0$. The channel $\mathcal N$ is either pure loss with transmissivity $0<\eta<1$, or a quantum-limited amplifier with gain $G>1$.

\begin{equation}
\begin{gathered}
E_F(\rho_{AB}^{\otimes k})=kE_F(\rho_{AB}),\\
E_C(\rho_{AB})=E_F(\rho_{AB})=E_F^{\mathrm G}(\rho_{AB}).
\end{gathered}
\label{eq:ser12-progress-3-1}
\end{equation}

Here $E_F^{\mathrm G}$ is the restriction to Gaussian pure-state decompositions. \sourcecite{ref:ser12-3}{Wil18}
\end{itemize}

\subsection*{References}

\begin{enumerate}[label={},leftmargin=4.8em,labelsep=0.5em]
  \footnotesize
  \item[\textup{[YKHL25]}]\label{ref:ser12-1}
  H. Yamasaki, K. Kuroiwa, P. Hayden, and L. Lami, "Entanglement Cost for Infinite-Dimensional Physical Systems," \emph{Communications in Mathematical Physics} \textbf{406}, 277 (2025). \href{https://doi.org/10.1007/s00220-025-05431-1}{doi:10.1007/s00220-025-05431-1}; \href{https://arxiv.org/abs/2401.09554}{arXiv:2401.09554}.

  \item[\textup{[Ade26]}]\label{ref:ser12-2}
  G. Adesso, "Optimality of Gaussian Entanglement of Formation," preprint (2026), version 2, 13 August 2026. \href{https://doi.org/10.48550/arXiv.2608.01909}{doi:10.48550/arXiv.2608.01909}; \href{https://arxiv.org/abs/2608.01909v2}{arXiv:2608.01909v2}.

  \item[\textup{[Wil18]}]\label{ref:ser12-3}
  M. M. Wilde, "Entanglement Cost and Quantum Channel Simulation," \emph{Physical Review A} \textbf{98}, 042338 (2018). \href{https://doi.org/10.1103/PhysRevA.98.042338}{doi:10.1103/PhysRevA.98.042338}; \href{https://arxiv.org/abs/1807.11939}{arXiv:1807.11939}.
\end{enumerate}

\subsection*{Comment}

The missing inequality must exclude collective non-Gaussian decompositions of many copies. A single-copy covariance optimization does not provide this exclusion. The cost is measured in asymptotic Bell pairs per copy with vanishing trace-norm error under unrestricted LOCC.

\subsection*{Field}

Quantum Resource Theory

\subsection*{Topic}

Gaussian quantum information; Entanglement cost; Additivity and regularization

\subsection*{ID}

\texttt{op\_c1499a0073ee800c}
